%===============================================================================
%      
%===============================================================================
\section{Introduction}

   Les interruptions sont un outil fondamental dans la communication
d'un système avec son environnement. Leur gestion, nécessaire, soulève
quelques questions liées notemment au fil d'exécution.

%-------------------------------------------------------------------------------
%
%-------------------------------------------------------------------------------
\subsection{Interruption et exception}

   Une interruption peut être vue commme la manifestation d'un
événement qui va perturber le fil d'exécution en cours sur un
processeur. Nous pouvons les classifier en deux catégories

\begin{description}
   \item[Les interruptions asynchrones] (que nous appellerons
     simplement {\rm interruptions}) sont le fait d'un élément
     extérieur au fil d'exécution en cours (une touche du clavier a
     été manipulée, une carte réseau a reçue une trame, \ldots). Nous
     parlerons également d'{\sc irq} (ou {\em Interrupt ReQuest} pour
     demande d'interruption).
   \item[Les interruptions synchrones] (que nous appellerons {\em
     exceptions}) sont déclanchées par le processeur lui-même comme
     conséquence de l'exécution (ou de la tentative d'exécution) d'une
     instruction. 
\end{description}

   Une interruption, qu'elle soit synchrone ou non, aura pour
conséquence l'exécution par le processeur d'un code, appelé
{\em handler} d'interruption. Ce code est en charge de traiter
l'interruption.

%...............................................................................
%
%...............................................................................
\subsection{Le {\em Programmable Interrupt Controller}}
   Les interruptions peuvent arriver depuis plusieurs sources. Celles
provenant de périphériques (clavier, cartes réseau, \ldots) sont
éventuellement acheminées au travers de circuits spécialisés appelés
{\sc pic}.

   Le {\sc pic} est donc un circuit électronique permettant de gérer les
{\sc irq} provenant de l'environnement afin de les délivrer au
processeur (c'est en fait une fonction du {\em southbridge}).

   Je ne m'étalerai pas ici sur les différents types de {\sc pic} ni
sur leur fonctionnement. Ce qu'il est nécessaire de comprendre, c'est
qu'il existe donc un circuit avec lequel le processeur va devoir
dialoguer pour une partie de la gestion des interruptions.

   Nous verrons plus loin que ManuX gère le 8259A qui est un {\sc pic}
conçu par Intel.

%-------------------------------------------------------------------------------
%
%-------------------------------------------------------------------------------
\subsection{Masquage}

   Une interruption peut éventuellement être  {\em masquée}, ce qui
signifie qu'elle ne sera pas délivrée au processeur, qui ne
s'appercevra donc pas qu'un événement s'est produit. Mais le masquage
d'une interruption n'empêche évidemment pas cet événement d'arriver !

%-------------------------------------------------------------------------------
%
%-------------------------------------------------------------------------------
\subsection{Traitement d'une interruption : partie haute et partie basse}

    Le traitement d'une interruption est donc en partie réalisé dans
un {\em handler}. Pour des raisons de simplicité, nous allons écrire
ces handlers en assembleur mais nous y utiliserons dès que possible
des fonctions C.

   Il est fondamental de comprendre que lorsqu'une interruption est
prise en compte par le processeur, le code qu'il était en train
d'exécuté est {\em interrompu} (surprise ?) pour exécuter le {\em
  handler}. Lorsque l'exécution de ce dernier s'achève, l'exécution du
code interrompu reprend son cours. Il est donc crucial que le {\em
  handler} ne réalise pas de modification du système qui rendrait
cette exécution incohérente (par exemple en supprimant une ressource
que ce code s'appétait à utiliser après s'être assuré de sa
disponibilité).

   Un autre élément fondamental est le fait que lorsqu'un handler
s'exécute, les interruptions ont inhibées et ne peuvent donc plus être
prises en compte\footnote{C'est un peu plus compliqué sur un système
  multiprocesseur.}. 
    
%===============================================================================
%      
%===============================================================================
\section{Le {\sc pic} Intel 8259A}

   Sur un (vieux) {\sc pc}, le {\sc pic} est un circuit Intel 8259A. Ce
circuit permet de gérer 8 types d'interruptions et on en utilise
plusieurs en cascade pour les additionner. Nous supposerons ici qu'il
y en a deux, le {\sc pic} 2 étant derrière le {\sc pic} 1 (en
``cascade'') pour un total de 15 interruptions.

   Ce circuit, donc on peut trouver la spécification
\cite{intel-8259a-spec} se programme au travers de deux registres

\begin{description}
    \item[Le registre de commande] accessible par le port 0x20 pour le
      circuit maître et 0xa0 pour le circuit esclave ;
    \item[Le registre de données] accessible par le port 0x21 pour le
      circuit maître et 0xa1 pour le circuit esclave.
\end{description}

   Une phase d'initialisation est nécessaire afin de configurer
correctement ces circuits.
   
%...............................................................................
%
%...............................................................................
\subsection{Gestion du {\sc pic} Intel 8259a dans ManuX}

   Ce circuit est géré au travers de fonctions définies dans {\tt
include/manux/intel-8259a.h} et implantées dans {\tt
     i386/intel-8259a.c}.

   L'objectif est assez simple : fournir (aux drivers des matériels
utilisant un mécanisme d'interruption) une interface générique,
c'est-à-dire indépendante du {\sc pic} utilisé. Nous allons donc
définir quelques fonctions de base permettant à ces drivers de
s'interfacer avec le système d'interruption. Si on souhaite utiliser
un autre circuit {\sc pic} (par exemple un  {\sc apic}), il
``suffira'' de réécrire ces fonctions, il n'y aura pas besoin de
modifier le code des nombreux drivers (au moins deux, rendez-vous
compte du boulot !) qui l'utilisent.

   Le principe général consiste à définir un handler d'interruption
unique qui va prendre en charge toutes les interruptions générées au
travers du {\sc pic}. Ce handler traitera les aspects spécifiques au
{\sc pic} puis invoquera le handler de chacun des ``clients''
c'est-à-dire des drivers des matériels qui sont à la source des
interrruptions.

   Décrivons les fonctions principales. Vous retrouverez leur code
dans {\tt i386/intel-8259a.c}.

\paragraph{\tt i8259aInit} permet d'initialiser
les circuits (le {\sc pic}) et de les configurer pour que les
interruptions générées utilisent une plage de numéros commençant à une
valeur passée en paramètre.
\paragraph{\tt i8259aAjouterHandler} permet d'associer une fonction
et des données à une {\sc irq}. C'est par le biais de cette fonction
qu'un client va se faire connaître auprès de notre système.
\paragraph{\tt i8259aGestionIRQ}
   
%===============================================================================
%      
%===============================================================================
\section{Définition des handlers}

   Sur le 386, une table (l'{\sc idt} pour {\em Interrupt Descriptor Table})
définit le comportement à adopter pour chacune des 256 interruptions
et exceptionns. Cette table est enregistrée grâce à l'instruction
\lstinline!lidt!.

   Cette table contient des entrées de trois types différents

\begin{description}
   \item[Interrupt Gate]       
   \item[Trap Gate Descriptor]       
   \item[Task Gate Descriptor]       
\end{description}

%-------------------------------------------------------------------------------
%
%-------------------------------------------------------------------------------
\subsection{Gestion dans ManuX}

   Observons ici comment tout cela est mis en place dans ManuX.

   Nous ne pouvons pas ici faire l'impasse sur un peu d'assembleur, si
bien que les fonctions ``bas niveau'' seront implantées dans les
fichiers {\tt manux/i386/interBasNiveau.h
  i386/interBasNiveau.nasm}. Tout le reste sera en C dans {\tt
  manux/interruptions.h} et {\tt i386/interruptions.c}.

   C'est la fonction C \lstinline!initialiserIDT()! qui est invoquée
par \lstinline!main()! pour initialiser la table.
   
%-------------------------------------------------------------------------------
%      
%-------------------------------------------------------------------------------
\subsection{Définition ``bas niveau'' d'un handler}

   Voici par exemple comment est écrit un handler pour l'interruption
7 (mais c'est sans intérêt ici) dont le rôle est d'invoquer une
fonction C nommée \lstinline!handlerPanique!. Notez que ce handler
empile tous les registres ainsi que le numéro de l'interruption
concernée. 

\begin{lstlisting}{language=nasm}
; Un handler qui va afficher un message
;--------------------------------------
stubHandlerPanique 7
  
         pusha
	 push dword 07h

         call handlerPanique
	 
         add esp, 4
	 popa

         iret
\end{lstlisting}

   Vous trouverez également un handler qui ne fait rien :

\begin{lstlisting}{language=nasm}
; Un handler qui ne fait rien ...
;--------------------------------
stubHandlerNop :
        iret                        ; On revient ...
\end{lstlisting}

   Notons que lors d'une interruption, le microprocesseur empile les
valeurs de {\tt EFLAGS}, {\tt CS} et {\tt EIP} (voir par exemple
\cite{intel386pru1986} page 159).

   Voyons maintenant comment écrire la suite plus confortablement en
C.
   
%-------------------------------------------------------------------------------
%      
%-------------------------------------------------------------------------------
\subsection{Définition d'un handler en C}

   Nous pouvons par exemple écrire la fonction
\lstinline!handlerPanique! dont le rôle est d'afficher un message
permettant de déterminer la cause de l'interruption :

\begin{lstlisting}
void handlerPanique(uint32_t itNum, TousRegistres registres,
		    uint32_t eip, uint32_t cs, uint32_t eFlags)
{
   ...
}
\end{lstlisting}

   Nous la définissons avec les paramètres qui sont, dans l'ordre
inverse, les valeurs empilées par le processeur lors de l'interruption
puis par le handler ``bas niveau'' décrit ci-desssus.

   La structure \lstinline!TousRegistres!, décrite dans {\tt
manux/types.h}, permet de manipuler simplement l'ensemble des
registres.
   
%-------------------------------------------------------------------------------
%      
%-------------------------------------------------------------------------------
\subsection{Insertion d'un handler dans l'{\sc idt}}

   La procédure suivante initialise une entrée dans l'{\sc idt} :

\begin{lstlisting}{}
void positionnerHandlerInterruption(IDT idt, int i, Handler handler)
/*
 * Affectation du handler de l'interruption i
 */
{
   idt[i].itg.offsetFaible = ((uint32_t)handler & 0xFFFF);
   idt[i].itg.offsetFort = ((uint32_t)handler >> 16);
   idt[i].itg.parametres = 0x8E00;
   idt[i].itg.selSegment = SELECTEUR_SEGMENT_CODE;
}

\end{lstlisting}

   Elle sera donc appelée par la procédure d'initialisation de l'{\sc idt} pour
chacune des interruptions.

%-------------------------------------------------------------------------------
%      Initialisation de la table
%-------------------------------------------------------------------------------
\subsection{Initialisation de la table}

   Nous avons maintenant tous les éléments pour construire notre {\sc
idt}. C'est le rôle de la fonction C \lstinline!initialiserIDT()!.
   
   Nous y positionnons donc tous les handlers sur le handler minimaliste
qui ne fait rien :

\begin{lstlisting}{}
   /* Comportement par defaut : on ne fait rien ! */
   for (i = 0; i < 256; i++) {
      positionnerHandlerInterruption(idt, i, stubHandlerNop);
   }
\end{lstlisting}

   Nous allons ensuite définir un handler un peu plus utile pour les
interruptions que peut nous générer le matériel

\begin{lstlisting}{}
   positionnerHandlerInterruption(idt, 0, stubHandlerPanique_0);
   positionnerHandlerInterruption(idt, 1, stubHandlerPanique_1);
   ...
\end{lstlisting}
   
   Nous pouvons enfin définir les handlers plus spécifiques

\begin{lstlisting}{}
   /* Le handler du timer */
   positionnerHandlerInterruption(idt, intTimer, stubHandlerTimer);

   /* Le handler du clavier */
   positionnerHandlerInterruption(idt, intClavier, stubHandlerClavier);
\end{lstlisting}


   
